\section{Method}
\label{sec:method}

A page is rendered once with a headless browser (Appendix~\ref{sec:appendix-figures}, Figure~\ref{fig:pipeline} diagrams the flow); its DOM geometry drives an adaptive tiling step (\S\ref{sec:tiling}); tiles are organized into a two-level graph, optionally spanning several hyperlinked pages (\S\ref{sec:graph}); a VLM reader is contrastively fine-tuned to score (question, tile) pairs (\S\ref{sec:contrastive}); and, at query time, an online controller walks the graph under a budget to decide which nodes the reader actually opens (\S\ref{sec:controller}).

\subsection{DOM-Guided Adaptive Tiling}
\label{sec:tiling}

\paragraph{Rendering and extraction.} We load the page in headless Chromium at a fixed viewport width (2048\,px), remove interface chrome via a CSS selector list, and take a full-page screenshot. For a configurable tag set (\code{h1}--\code{h4}, \code{p}, \code{table}, \code{figure}, \code{img}, \code{ul}/\code{ol}, \code{blockquote}, \code{pre}) we read each element's absolute bounding box from the rendered DOM. For non-text containers this is \code{getBoundingClientRect}; for text-flow tags, whose CSS box spans the full container width regardless of actual content, we instead union the client rects of every descendant text node. Without this, a paragraph or reference list sharing a row with a floated infobox reports a bounding box running through the infobox's column even though its text stops well short of it. Hyperlinks are collected from \code{<p>} elements only, seeding the cross-page edges of \S\ref{sec:graph}.

\paragraph{Overlap resolution.} Nested tags (an \code{<img>} in a \code{<figure>}) yield duplicate bounding boxes for the same content; we sort by DOM depth and drop any element already fully contained (1\,px tolerance) in a shallower kept element. A floated sidebar and the text sharing its row can also genuinely intersect even after the text-flow tightening above. Rather than tile such elements independently and duplicate the shared pixels, we connect any two elements whose bounding-box intersection exceeds a small fraction of the smaller one's area, and tile each connected component as one crop spanning their union.

\paragraph{Adaptive tiling and patching.} Elements are processed top-to-bottom under three rules: (1) \emph{heading pairing} -- a heading is never a standalone tile, but merged with the next element (e.g.\ \code{h2+p}); (2) \emph{merging} -- non-heading elements under 80\,px are buffered until their union spans that minimum, then flushed as one \code{merged} tile; (3) \emph{patching} -- an element taller than 2048\,px is sliced into $\lceil h/h_{\max}\rceil$ equal-height sub-tiles. This targets PixelRAG's failure mode directly: a table that would straddle a fixed grid cell is instead one tile, and a short caption merges with its neighbor rather than being orphaned. Each tile keeps a pointer to its source element(s), used both for the graph (\S\ref{sec:graph}) and as ground truth for negative-mining (\S\ref{sec:contrastive}). Figure~\ref{fig:tiling-overlay} (Appendix~\ref{sec:appendix-figures}) shows tiles overlaid on a real page.

\subsection{Multigranular Graph Construction}
\label{sec:graph}

Tiles are organized into a two-level directed graph (one \code{page} node, one \code{element} node per tile) with all five of MAGE-RAG's relation types, plus one of our own for cross-page connectivity:
\begin{itemize}
\item \textbf{\code{contains}} ($\mathrm{page}\to\mathrm{element}$): every element belongs to its page.
\item \textbf{\code{reading\_order}} ($\mathrm{element}_i\to\mathrm{element}_{i+1}$): a spatial heuristic, \emph{not} DOM order -- elements are grouped into visual rows (20\,px vertical-overlap tolerance), each sorted left-to-right, rows concatenated top-to-bottom, since DOM and visual order diverge under CSS positioning.
\item \textbf{\code{layout\_adjacency}} ($\mathrm{element}\leftrightarrow\mathrm{element}$): reciprocal edges to the 2D neighbors reading order flattens away (left/right within a row; above/below to the nearest horizontally-overlapping tile in the next row).
\item \textbf{\code{section\_hierarchy}} ($\mathrm{heading}\to\mathrm{child}$): each heading points to its nested sub-heading/content, forming a section tree distinct from the flat \code{contains} relation.
\item \textbf{\code{semantic\_neighbor}} ($\mathrm{element}\leftrightarrow\mathrm{element}$): TF-IDF-cosine edges between lexically close elements regardless of position, added greedily under a per-tile degree cap. Optional, off by default, meant to be toggled to compare the controller with and without it (\S\ref{sec:experiments}).
\item \textbf{\code{links\_to}} ($\mathrm{element}\to\mathrm{external\_page}$): hyperlinks from the first $k{=}3$ link-bearing text tiles (reading order); beyond that budget, links are dropped to bound the corpus crawl's branching factor.
\end{itemize}
Every element node starts with a \code{state} attribute (\state{inactive}), which the controller (\S\ref{sec:controller}) evolves at query time. Figure~\ref{fig:page-graph} (Appendix~\ref{sec:appendix-figures}) shows the resulting graph for the page of Figure~\ref{fig:tiling-overlay}.

\paragraph{Oversized-page splitting.} One article's screenshot exceeded $46{,}000$\,px and produced dozens of tiles whose sequential QA generation (\S\ref{sec:contrastive}) ran for hours with no partial progress saved on failure. Rather than exclude such pages, one whose element height exceeds $10\times$ the max tile height is split into sections \emph{before} tiling, at \code{h1}/\code{h2} boundaries, with any still-oversized section further re-chunked by a height budget. Each section is tiled and graphed as a standalone page would be; the per-section graphs are merged by node-namespacing, with consecutive sections' \code{page} nodes linked by a new \textbf{\code{continues}} relation -- distinct from \code{links\_to}: sequential parts of the same document, not a hyperlink elsewhere. Every page's graph, split or not, exposes an \emph{entry point} for the corpus graph below to rewire onto; the controller does not yet cross \code{continues} edges automatically (\S\ref{sec:limitations}). Figure~\ref{fig:section-split} (Appendix~\ref{sec:appendix-figures}) shows an example split into 16 sections.

\paragraph{Corpus graph.} From a seed page we crawl up to $m$ pages reachable via its \code{links\_to} edges (one hop, non-recursive), build each linked page's graph independently, and merge by namespacing every node with a URL-derived slug, rewiring each placeholder external-page node onto the crawled page's entry point. \code{links\_to} is the only cross-document connective tissue, \code{continues} the only cross-section one -- the substrate the evidence controller (\S\ref{sec:controller}) operates over.

\paragraph{Structured rendering.} Any (sub-)graph serializes its page node and directly-contained elements to XML (one \code{<element id=... tag=... state=...>text</element>} per node) as the reader's structured input, in the spirit of MAGE-RAG's structured-rendering step (example in Appendix~\ref{sec:appendix-figures}, Figure~\ref{fig:xml}).

\subsection{Contrastive Reader Fine-Tuning}
\label{sec:contrastive}

To turn a general-purpose VLM into a tile-level reader that also scores how well a tile answers a query, we fine-tune it contrastively on synthetic data, in three stages.

\paragraph{Synthetic QA generation.} For every tile, we prompt the VLM with the tile image alone for one self-contained, factual question answerable \emph{only} from that tile, or \code{SKIP} if it carries no exploitable content -- the recipe of synthetic query generation for retriever training (InPars, Promptagator;~\citealp{inpars2022,promptagator2022}) applied to image tiles. Every tile crop is persisted regardless of outcome, so \code{SKIP}ped tiles remain available as candidate hard negatives.

\paragraph{Hard-negative mining.} For each question, we rank every other tile from the same page by TF-IDF cosine similarity and keep the top $n{=}2$ clearing two false-negative filters: (a) a lexical filter rejecting any tile whose text already contains the answer verbatim, and (b), optionally, a VLM-judge filter showing the candidate image to the reader VLM and asking whether it can recover the answer there, catching non-lexical false negatives at the cost of one extra VLM call per candidate.

\paragraph{Contrastive LoRA fine-tuning.} \code{Qwen2-VL} is generative, not a similarity encoder; we turn it into one via PromptEOL~\citep{prompteol2024} (a short fixed instruction, ``summarize the above in one word,'' taking the last token's final-layer hidden state as embedding), originally proposed for text-only LLMs and applied here to a VLM's image+instruction input. LoRA ($r{=}16$, $\alpha{=}32$, dropout $0.05$) adapts both the decoder's attention projections and the vision tower's, so the model can refine what it attends to in the image jointly with how it reads the question. Given $B$ questions each with one positive and $K$ mined negatives, we minimize a temperature-scaled ($\tau{=}0.05$) InfoNCE loss over the $B(1{+}K)$ tiles in the batch, so each positive competes against its own negatives and, in-batch, every other question's tiles.

\paragraph{Resolution as an efficiency lever.} PixelRAG reports up to $3\times$ token-cost reduction at lower resolution without specifying the filter or resolutions tested. We build the infrastructure to test this directly: every tile can be mirrored via Lanczos resampling at a configurable scale (0.5 default), reusing the same questions/positives/negatives at both resolutions so comparing them is a matter of pointing fine-tuning at one image root or the other (\S\ref{sec:experiments}).

\subsection{Online Evidence Controller}
\label{sec:controller}

Given a query, the controller decides which element nodes the reader actually reads, rather than the whole page or a static top-$k$. It evolves each node's \code{state} through the four-state machine of Figure~\ref{fig:states}, generalizing MAGE-RAG's Algorithm~1.

\begin{figure}[t]
\centering
\begin{tikzpicture}[
  scale=0.8, transform shape,
  node distance=1.1cm and 1.6cm,
  st/.style={draw, rounded corners, minimum width=1.9cm, minimum height=0.65cm, align=center, font=\small},
  arr/.style={-{Stealth[length=2mm]}, thick}
]
\node[st] (inactive) {\state{inactive}};
\node[st, right=of inactive] (active) {\state{active}};
\node[st, above right=0.5cm and 1.5cm of active] (opened) {\state{opened}};
\node[st, below right=0.5cm and 1.5cm of active] (pruned) {\state{pruned}};
\draw[arr] (inactive) -- node[above, font=\scriptsize]{activate} (active);
\draw[arr] (active) -- node[above, sloped, font=\scriptsize]{open (budget --1)} (opened);
\draw[arr] (active) -- node[below, sloped, font=\scriptsize]{prune (free)} (pruned);
\end{tikzpicture}
\caption{The controller's state machine. Opening a node consumes budget and may activate inactive \code{reading\_order} neighbors; pruning is free, applied to any node still \state{active} once budget is exhausted.}
\label{fig:states}
\end{figure}

\paragraph{Relevance scoring.} Every element node is scored against the query by TF-IDF cosine similarity (same mechanism as hard-negative mining), a deliberate placeholder for the fine-tuned reader's own learned similarity (\S\ref{sec:limitations}).

\paragraph{Policy.} Algorithm~\ref{alg:controller}: (1) seed the frontier with the $k$ highest-scoring inactive nodes; (2) repeatedly take the highest-scoring \emph{active} node, pruning it for free below an opening threshold or opening it (one unit of budget $B$) and activating still-inactive \code{reading\_order} neighbors that clear an activation threshold; (3) once budget is exhausted or no active node remains, prune whatever is left active. Neighbors are only considered \emph{after} their node is opened, so the frontier grows along reading order conditioned on what has already proven relevant, rather than a fixed global ranking -- what distinguishes it from static top-$k$ retrieval over the same scores (ablation (d) in \S\ref{sec:experiments}).

\begin{algorithm}[t]
\caption{Budgeted best-first evidence control}
\label{alg:controller}
\begin{algorithmic}[1]
\Require graph $G$, query $q$, budget $B$, seed count $k$, thresholds $\theta_{\text{open}}, \theta_{\text{act}}$
\State score every element node by $\mathrm{sim}(q, \text{text}(n))$
\State activate the $k$ highest-scoring inactive nodes \Comment{seed}
\State $b \gets 0$
\While{$b < B$ and some node is \state{active}}
  \State $n \gets \arg\max$ score over \state{active} nodes
  \If{$\mathrm{score}(n) < \theta_{\text{open}}$}
    \State $n \gets$ \state{pruned} \Comment{free}
  \Else
    \State $n \gets$ \state{opened}; $b \gets b + 1$; record its text as evidence
    \For{each inactive \code{reading\_order} neighbor $n'$ of $n$}
      \If{$\mathrm{score}(n') \geq \theta_{\text{act}}$}
        \State $n' \gets$ \state{active}
      \EndIf
    \EndFor
  \EndIf
\EndWhile
\State prune every node still \state{active} \Comment{finalize}
\State \Return concatenation of \state{opened} nodes' text, in reading order
\end{algorithmic}
\end{algorithm}

\paragraph{Output.} The concatenation of \state{opened} nodes' text, in reading order, is the final evidence passed to the downstream reader, with the controller's action log available for inspection.
